\begin{figure}[t]
    \centering
    \begin{subfigure}[b]{\textwidth}
        \centering
        \resizebox{\linewidth}{!}{
        \begin{tikzpicture}
        \begin{axis}[
            xlabel={Number of clusters},
            ylabel={\gls{cu}},
            grid=major,
            xmin=1, xmax=20,
            ymin=0, ymax=1,
            ytick distance=0.25,
            ymajorgrids=true,
            width=\linewidth,
            height=0.45\linewidth,
        ]
            \addplot[color=breakoutColor, mark=\breakoutMark, mark size=\breakoutMarkSize] coordinates {
                (2,0.66) (3,0.57) (4,0.65) (5,0.79) (6,0.59) (7,0.62) (8,0.56) (9,0.56) (10,0.58) (11,0.64) (12,0.59) (13,0.55) (14,0.55) (15,0.54) (16,0.59) (17,0.60) (18,0.53) (19,0.55) (20,0.55)
            };
        \end{axis}
        \end{tikzpicture}
        }
        \caption{\gls{cu} with respect to the number of clusters.}
    \end{subfigure}

    \vspace{1em}

    \begin{subfigure}[b]{\textwidth}
        \centering
        \resizebox{\linewidth}{!}{
        \begin{tikzpicture}
        \begin{axis}[
            xlabel={Number of clusters},
            ylabel={\gls{as}},
            grid=major,
            xmin=1, xmax=20,
            ymin=0, ymax=1,
            ytick distance=0.25,
            ymajorgrids=true,
            width=\linewidth,
            height=0.45\linewidth,
        ]
            \addplot[color=breakoutColor, mark=\breakoutMark, mark size=\breakoutMarkSize] coordinates {
                (2,0.62) (3,0.83) (4,0.74) (5,0.82) (6,0.85) (7,0.85) (8,0.85) (9,0.85) (10,0.83) (11,0.82) (12,0.85) (13,0.85) (14,0.84) (15,0.83) (16,0.83) (17,0.83) (18,0.84) (19,0.85) (20,0.83)
            };
        \end{axis}
        \end{tikzpicture}
        }
        \caption{\gls{as} with respect to the number of clusters.}
    \end{subfigure}

    \vspace{1em}

    \begin{subfigure}[b]{\textwidth}
        \centering
        \resizebox{\linewidth}{!}{
        \begin{tikzpicture}
        \begin{axis}[
            xlabel={Number of clusters},
            ylabel={\gls{acu}},
            grid=major,
            xmin=1, xmax=20,
            ymin=-0.2, ymax=1,
            ytick distance=0.2,
            ymajorgrids=true,
            width=\linewidth,
            height=0.45\linewidth,
        ]
            \addplot[color=breakoutColor, mark=\breakoutMark, mark size=\breakoutMarkSize] coordinates {
                (2,0.28) (3,0.40) (4,0.39) (5,0.61) (6,0.43) (7,0.47) (8,0.41) (9,0.41) (10,0.42) (11,0.46) (12,0.45) (13,0.40) (14,0.39) (15,0.37) (16,0.42) (17,0.43) (18,0.36) (19,0.40) (20,0.38)
            };
        \end{axis}
        \end{tikzpicture}
        }
        \caption{\gls{acu} with respect to the number of clusters.}
    \end{subfigure}
    \caption{\gls{cu}, \gls{as}, and \gls{acu} with respect to the number of clusters, for layer 10 of the Breakout surrogate policies (the layer with the highest average \gls{acu}).}
    \label{fig:latent_space_curves_breakout}
\end{figure}
